% Owner: classification_e. Root and boundary_d full-source referee PASS.
\begin{proposition}[E.2.21 (Proved): global identification of a smooth completion]
Let $\pi:\mathcal X\to\Delta$ satisfy the hypotheses of
E.2.5, with fixed coprime multiplicities $(3k,4l)$ on
every noncentral fibre.  Then $X_0$ is a member of that
same constructed family.  Its isomorphism parameter is
the interior value $z(0)$ of E.2.7.

After choosing a holomorphic lift $c(s)$ of that parameter
and shrinking the disc, the actual family is isomorphic
to the prescribed family at $c(s)$.  The ambiguity in
this comparison is the connected translation automorphism
group already determined in E.1.
\end{proposition}
\begin{proof}
E.2.10 extends the intrinsic pencil to a proper flat map
\[
 g:\mathcal X\longrightarrow\mathbb P^1\times\Delta.
\]
E.2.7 supplies its interior nonzero period $z(s)$ and hence
a holomorphic logarithmic lift $c(s)$.  E.2.12 excludes
additional critical values, and E.2.14--E.2.15 exclude
target collapse.  Normalize the three distinct labelled
sections to $0,1,\infty$ by a holomorphic projective
transformation over the disc.

On the ordinary locus the actual marked Jacobian has
period matrix $\Pi_{c(s)}(t)$.  More intrinsically, its
holomorphic period kernel in the complexified marked
lattice equals that of this matrix.  This equality holds
for $s\ne0$ by the given family identification and extends
over $s=0$ by the identity theorem.  The local marking
choices agree: an automorphism of a fixed marked member
has identity action on ordinary $H_1$, by E.1.2.  The
actual smooth family at any fixed ordinary target point
also has identity homology return around $s=0$.

E.2.13 gives smooth actual normalized root covers at the
two finite values, of degrees $3k$ and $4l$, with their
free cyclic deck actions.  Their marked Jacobians are
the corresponding reference root Jacobians.  E.2.20
identifies the actual cusp neighbourhood over the full
$(q,s)$-bidisc with the prescribed marked Mumford
neighbourhood, including its period multipliers.
It remains to identify the finite affine actions and
the global translations on the overlaps.

\emph{The finite affine classes.}
Choose local holomorphic origins on a finite root torus
family.  Every deck map is affine on the compact torus
fibres.  Its linear part is the fixed marked finite
monodromy.  Subtracting this part gives a holomorphic
translation section.  Lift it to the Lie bundle on
simply connected refinements of the root charts.
The finite-order relation says that the sum of its
successive linear translates is an integral period.
This lattice vector is locally constant in $s$:
the lifted translation and period vectors vary
continuously, whereas the integral coefficients are
discrete.  Changing the lift changes it by precisely
the integral coboundary allowed in E.1.

Consequently the finite affine cocycle has the same
integral class at $s=0$ as at $s\ne0$.  After matching
these classes, the remaining difference is a cocycle
in the complex vector sheaf.  The explicit finite
averaging in E.1 removes it, since division by the
finite group order is allowed in that sheaf.
This constructs the finite-root comparisons with
holomorphic dependence on $s$.  Equivalently one
may use a finite chart cover and its deck translates;
the same averaging solves the equivariant comparison.
The reduced finite fibres are the resulting free
quotients; they have not been assumed to be tori.

\emph{The integral collar classes.}
Use the two finite-root neighbourhoods, the cusp
neighbourhood, ordinary simply connected charts,
and fixed annular collars to cover the target.
Choose local holomorphic origins on the ordinary
torus charts.  A transition between two such
descriptions is fibrewise affine.  The common
integral marking makes its linear part the
prescribed one, leaving a translation section.

Lift these sections to the Lie bundles on simply
connected refinements of the collars.  Differences
on deck translates and sums on triple overlaps
are integral periods.  Every resulting integral
vector is constant in $s$.  This comparison takes
place on fixed collars away from the singular
centres; it does not evaluate a logarithm at a
singular fibre.  The finite and cusp chart
identifications specify exactly which integral
changes of origins are allowed.

For $s\ne0$ these classes are those of the
prescribed logarithmic-transform data, modulo
the integral changes of origins in E.1.
They are therefore the same at $s=0$.
One can normalize them once on the connected
parameter disc.  This includes the cusp
winding class, since E.2.20 identifies the
actual marked period action there.  Thus
there is no additional central integral
translation or fan datum.

\emph{Solving the holomorphic remainder.}
After this discrete normalization, the remaining
comparison is the additive cocycle in the
extension sheaf $\mathcal V$ of E.1.  Its
finite extension conditions are exactly
those of the already identified root
charts.  Its cusp extension condition
is the actual toric collar condition
proved in E.2.20.  The calculation in
E.1 therefore applies to this actual
cover and gives
\[
 0\longrightarrow\mathcal O_{\mathbb P^1}
       \boldsymbol\delta
 \longrightarrow\mathcal V
 \longrightarrow\mathcal O_{\mathbb P^1}(-\infty)
 \longrightarrow0,
 \qquad H^1(\mathbb P^1,\mathcal V)=0.
\]
The explicit finite averaging and Laurent-series
splitting in E.1 kill the additive cocycle.
The resulting translations glue the local
comparisons to a biholomorphism of $X_0$
with the prescribed member at $c(0)$.

The same solution works over a smaller parameter
disc.  On the fixed collars, Cauchy integral
formulas for the Laurent coefficients preserve
holomorphic dependence on $s$, as does finite
averaging.  Thus the solved translations and
their glued comparison are holomorphic in
$s$.  Their ambiguity is the connected
scalar translation group of E.1.  This
proves both the central and family assertions.
\end{proof}

\begin{remark}[Scope of E.2.21]
The known-pair locus is open by E.2.2 and
is now stable under full punctured-disc
smooth specializations.  These two facts
alone do not prove that an arbitrary open
known locus in a reduced analytic base
is topologically closed.  A holomorphic
disc through its boundary need not have
its entire punctured image in that locus.
The needed generic propagation argument
is proved in E.2.22.  E.2.23--E.2.24 then
establish the arbitrary-base and continuous-family
conclusions asserted in E.2.4.
\end{remark}
